\section{Introduction}
\label{sec:tpc30-introduction}

The ultra-long complement of a truncated divisor model is not made
small merely by reflecting every long divisor across its target
hyperbola.  Reflection replaces a long divisor by a shorter
complement and makes the orbit fibers visible, but the resulting
family still contains enough short fibers to recover the natural
scale.  The preceding paper therefore used two kinds of shared
content in the reflected variables: a large gcd on long fibers and a
small primitive lcm on sparse fibers \citep{WangTPC29}.  Those
arguments are absolute and leave the small-reflected-content core.

This paper observes that the chosen divisor variables do not contain
all of the available common-divisor information.  For two physical
rows $m\ne n$, put
\begin{equation}
 N_m(j)=mj+h,
 \qquad
 N_n(j)=nj+h,
 \qquad
 \cG_{m,n}(j)=(N_m(j),N_n(j)).
 \label{eq:tpc30-intro-target-gcd}
\end{equation}
The full target gcd is independent of how either target is factored.
On the primitive orbit it satisfies
\begin{equation}
 \boxed{
 \cG_{m,n}(j)=(N_m(j),m-n),
 \qquad
 \cG_{m,n}(j)\mid m-n.}
 \label{eq:tpc30-intro-gcd-identity}
\end{equation}
The divisibility statement appeared in the earlier target-gcd
geometry for prime source rows \citep{WangTPC18}.  The new point here
is a weighted large-content pruning theorem for the actual opened rows
$m=\ell d$, the orbit variable, and all three raw ultra channels.

Let $Q$ be the row scale and $J$ the orbit length, with
\begin{equation}
 QJ\asymp X.
 \label{eq:tpc30-intro-QJ}
\end{equation}
Let $\cI$ be the physical orbit interval, of length
$O_{\mathscr D}(J)$.  For a fixed off-diagonal row pair, the exact values of
$\cG_{m,n}(j)$ belong to the divisor set of $|m-n|$.  If
$\cG_{m,n}(j)=c$, the orbit variable lies in one residue class modulo
$c$.  It follows that
\begin{equation}
 \#\{j\in\cI:(j,h)=1,\ \cG_{m,n}(j)>Z\}
 \ll_{\eps,\mathscr D}X^\eps
 \left(1+\frac JZ\right).
 \label{eq:tpc30-intro-fixed-row-occupancy}
\end{equation}
The physical row coefficients have one-row $\ell^1$ mass $O(Q)$.
Summing \eqref{eq:tpc30-intro-fixed-row-occupancy} over the row pairs
gives the main estimate.

\begin{theorem}[Common-target-content pruning, informal form]
\label{thm:tpc30-intro-main}
For each of the two mixed raw channels and the both-ultra raw channel,
the part on which $\cG_{m,n}(j)>Z$ satisfies
\begin{equation}
 \boxed{
 |\cS_{\cG>Z}|
 \ll_{\eps,h,\mathscr D}X^\eps
 \left(
 Q^2+\frac{JQ^2}{Z}
 \right).}
 \label{eq:tpc30-intro-main-bound}
\end{equation}
The same estimate holds for their sum.
\end{theorem}

Since $JQ^2\asymp XQ$ is the natural absolute scale, the two terms on
the right of \eqref{eq:tpc30-intro-main-bound} save respectively
$J^{-1}$ and $Z^{-1}$.  Thus if
\begin{equation}
 Q=X^{\beta+o(1)},\qquad
 J=X^{1-\beta+o(1)},\qquad
 Z\ge X^\eta,
 \label{eq:tpc30-intro-exponents}
\end{equation}
then every fixed
\begin{equation}
 \eta_0<\min\{\eta,1-\beta\}
 \label{eq:tpc30-intro-saving}
\end{equation}
is an available power saving.  No cancellation between M\"obius
signs or between distinct rows is used.

The factorization consequence is especially useful.  In the
both-ultra channel write
\begin{equation}
 N_m=rU,
 \qquad
 N_n=sV.
 \label{eq:tpc30-intro-factorizations}
\end{equation}
Put
\begin{equation}
 g=(r,s),\quad e=(U,V),\quad
 r=gr_0,\quad s=gs_0,\quad U=eU_0,\quad V=eV_0.
 \label{eq:tpc30-intro-primitive-factors}
\end{equation}
Then $(r_0,s_0)=(U_0,V_0)=1$ and the stronger exact identity is
\begin{equation}
 \boxed{
 \cG_{m,n}(j)
 =ge(r_0,V_0)(s_0,U_0).}
 \label{eq:tpc30-intro-cross-content}
\end{equation}
No assumption such as $(g,e)=1$ is needed.  The two final factors are
the cross-contents left after the within-side gcds are removed.
Therefore \eqref{eq:tpc30-intro-main-bound} gives a power-saving upper
bound for every retained sector on which the product in
\eqref{eq:tpc30-intro-cross-content} exceeds $Z$.  The factor
$g=(r,s)$ is the reflected content used in TPC-29;
$e=(U,V)$ and the two cross-contents are new to the present reduction's
bookkeeping.  In particular, the
present estimate remains effective when $(r,s)=1$.

At the exact high-$\beta$ sample inherited from TPC-28,
\begin{equation}
 \beta=\frac{267}{400},
 \qquad
 1-\beta=\frac{133}{400},
 \label{eq:tpc30-intro-high-beta}
\end{equation}
consider the exponent cell
\begin{equation}
 r,s=X^{1/4+o(1)},\qquad
 (r,s)=X^{o(1)},\qquad
 U,V=X^{3/4+o(1)},\qquad
 (U,V)=X^{2/5+o(1)}.
 \label{eq:tpc30-intro-high-cell}
\end{equation}
Its reflected lcm has exponent $1/2$, strictly above the orbit
exponent $133/400$, and its primitive reflected lcm has the same
exponent.  Hence the TPC-29 absolute-incidence saving is zero on this
cell.  By contrast, the cross-content exponent is $2/5$, and the new
relaxed margin is
\begin{equation}
 \min\left\{\frac25,\frac{133}{400}\right\}
 =\boxed{\frac{133}{400}}.
 \label{eq:tpc30-intro-high-margin}
\end{equation}
This is an upper bound for a support-compatible formal sector, not an assertion that
the sector carries nonzero arithmetic mass.

The result does not solve the residual problem.  Once all
polynomially large target content is removed, the remaining core has
\begin{equation}
 \cG_{m,n}(j)\le X^\eta
 \label{eq:tpc30-intro-small-core}
\end{equation}
for every fixed threshold used in the pruning, with the fully coprime
case $\cG_{m,n}(j)=1$ as its sharp model.  An absolute gcd argument
has no further power saving there.  The next plausible interface is a
signed dispersion theorem across normalized determinant windows or
determinant residue classes.  Fixed-form logarithmically averaged
two-point results such as \citep{Tao2016} do not by themselves give a
uniform estimate for the polynomially growing family retained here.

\Cref{sec:tpc30-interface} records the physical packet and the three
raw channels.  \Cref{sec:tpc30-target-geometry} proves the exact gcd
and orbit-class identities.  \Cref{sec:tpc30-row-orbit} proves the
weighted pruning theorem.  \Cref{sec:tpc30-cross-content} transfers it
to cross-factor content, and \cref{sec:tpc30-high-beta} gives the
exact rational ledger.  \Cref{sec:tpc30-comparison} separates the new
region from earlier content and near-row reductions and formulates
the next gate.  \Cref{sec:tpc30-certificate} records the finite exact
certificate and the claim boundary.
